% ============================================================
% TOUCH — Section 1.0: The Body as Signifying Practice
% ============================================================
% Kobena Mercer — "Black Hair/Style Politics": hair as
% signifying practice. Aesthetics and politics inseparable.
% The HairSynth makes this signifying practice literally
% generative: cultural meaning in hair becomes computational
% input. The sign generates sound.
% Hortense Spillers — Mama's Baby, Papa's Maybe: body vs.
% flesh. The flesh is the site of racial inscription.
% The uninvited touch is a claim on the flesh.
% The HairSynth reclaims the flesh as instrument on its
% own terms — generating knowledge and sound, not absorbing
% the claims of others.
% Political context: CROWN Act, workplace discrimination,
% school dress codes. The institutional machinery that
% polices Black hair as a site of control.
%
% Key texts:
%   Mercer — "Black Hair/Style Politics"
%   Spillers — "Mama's Baby, Papa's Maybe: An American
%   Grammar Book"
% ============================================================

\section*{Section 1.0: The Body as Signifying Practice}

\ifdraft
\lipsum[1-3]
\fi
